\scanpage{10}\mspage{9}
\makeatletter\tagsleft@true\makeatother
de \uline{compatibilité} dans $\unsure{\mathcal{M}|M}$ est une \deleted{\unsure{notion}} relation importante, qui se \unsure{surrajoute|surajoute} aux relations $\underset{\mathcal{M}}{\leq}$, $\underset{\mathcal{M}}{\ll}$. Moralement, pour récupérer l'atelier (à iso unique près) il suffit de connaitre sur $\mathcal{M}$ ces deux relations d'ordre, et la relation de compatibilité.
% slanted
\margin{au moins dans le cas modéré}

Les relations \deleted{\unsure{suivantes|restantes}} \inserted{(compatibilité, disjonction)} sont si importantes, qu'elles méritent chacune une notation lapidaire. Je propose
\[
\deleted{F\parallel G}\qquad F\parallel G \tag*{(1.14)}
\]
pour la disjonction, et
\[
\deleted{F\unsure{\mathrel{\tikz[baseline=0pt,x=1em,y=1em]{\draw (-.38,.62)--(.38,.25)--(-.43,-.10); \draw (.34,.62)--(-.32,.25)--(.38,-.10);}}|\mathrel{\tikz[baseline=0pt,x=1em,y=1em]{\draw (-.44,.26)--(0,.58)--(.44,.26)--(0,-.06)--cycle; \draw (0,.58)--(0,-.06);}}}G}\qquad F\mathrel{\tikz[baseline=0pt,x=1em,y=1em]{\draw (-.44,.26)--(0,.58)--(.44,.26)--(0,-.06)--cycle; \draw (0,.58)--(0,-.06);}} G\quad\text{ou}\quad
\tikz[baseline=(fg.base)]{\node[inner sep=1pt] (fg) {$F\mathrel{\tikz[baseline=0pt,x=1em,y=1em]{\draw (-.38,.62)--(.38,.25)--(-.43,-.10); \draw (.34,.62)--(-.32,.25)--(.38,-.10);}} G$}; \draw (fg.center) ellipse [x radius=2.12em,y radius=1.04em];}
\tag*{(1.15)}
\]
pour la compatibilité (NB $\mathrel{\tikz[baseline=0pt,x=1em,y=1em]{\draw (-.38,.62)--(.38,.25)--(-.43,-.10); \draw (.34,.62)--(-.32,.25)--(.38,-.10);}}$ est plus rapide \unsure{:|à} écrire que $\mathrel{\tikz[baseline=0pt,x=1em,y=1em]{\draw (-.44,.26)--(0,.58)--(.44,.26)--(0,-.06)--cycle; \draw (0,.58)--(0,-.06);}}$, $\unsure{\mathrel{\tikz[baseline=0pt,x=1em,y=1em]{\draw (-.10,.58)--(-.52,.22)--(-.10,-.08); \draw (.10,.58)--(.53,.22)--(.15,-.08); \draw (-.49,.21)..controls (-.22,.10) and (.10,.10)..(.48,.24); \draw (-.10,.39)--(-.29,.19)--(-.11,.02); \draw (.18,.39)--(.39,.24)--(.19,.05);}}|\deleted{\mathrel{\tikz[baseline=0pt,x=1em,y=1em]{\draw (-.10,.58)--(-.52,.22)--(-.10,-.08); \draw (.10,.58)--(.53,.22)--(.15,-.08); \draw (-.49,.21)..controls (-.22,.10) and (.10,.10)..(.48,.24); \draw (-.10,.39)--(-.29,.19)--(-.11,.02); \draw (.18,.39)--(.39,.24)--(.19,.05);}}}}$ mais est moins suggestif de l'idée combinatoire exprimée.

On n'a rien dit pour l'instant au sujet de la relation $\deleted{\ll}\preccurlyeq$ de raffinement sur $\unsure{\mathcal{F}|\mathfrak{F}}$. Elle est formellement de nature plus délicate, et est la notion la plus spécifique dans l'``Algèbre des figures''. La signification géométrique s'éclaire par les exemples, et les axiomes que \unsure{ceux-|ceux}%
